\chapter{Estimasi, Interval Kepercayaan, dan Uji Hipotesis}

Sebuah tim menguji pendeteksi serangan pada $50$ contoh berlabel, dan pendeteksi itu
menangkap $33$ di antaranya. Laporan yang ditulis malam itu berbunyi: recall $66$
persen. Pendeteksi kedua, diuji pada $100$ contoh, menangkap $53$: recall $53$ persen.
Rapat pagi memutuskan memakai yang pertama, karena $66$ lebih besar daripada $53$, dan
keputusan itu selesai dalam satu menit.

Aritmetikanya benar: $33/50=0{,}66$. Yang keliru adalah anggapan bahwa angka itu
menunjuk satu bilangan. Bila lima puluh contoh itu diganti lima puluh contoh lain dari
kumpulan yang sama, hasilnya bukan $33$ lagi; ia bergoyang, dan goyangannya tidak
kecil. Sesudah dua jam menulis program, tim itu sebenarnya memegang dua kalimat yang
berbeda jauh kekuatannya. Kalimat pertama, ``pada data ini terlihat $66$ persen'',
hampir tidak berisi apa pun. Kalimat kedua, ``laju tangkapan sebenarnya ada di suatu
tempat dalam selang ini'', berisi sesuatu yang bisa dipakai, dengan syarat kita tahu
seberapa lebar selang itu dan apa arti lebar tersebut.

Bab 19 sudah memberi dua alat: hukum bilangan besar, yang mengatakan rata-rata sampel
memusat ke nilai sebenarnya, dan teorema limit pusat, yang mengatakan galatnya menyebar
normal dengan simpangan baku $\sigma/\sqrt n$. Keduanya pernyataan tentang banyak
percobaan sekaligus, dan tidak satu pun menjawab pertanyaan rapat pagi tadi: dari satu
percobaan yang sudah selesai, kalimat apa yang boleh diucapkan tentang parameter yang
tidak pernah dilihat, dan seberapa sering kalimat itu salah bila prosedurnya diulang.
Bab ini menyusun jawabannya, dari penaksir sampai nilai-p.

\begin{untukanak}
Seorang petani ingin tahu berapa banyak rumput liar di sawahnya. Seluruh sawah terlalu
luas untuk diperiksa, jadi ia mengambil dua puluh petak kecil yang tersebar, lalu
menghitung rumput di setiap petak. Dari dua puluh angka itu ia menebak keadaan seluruh
sawah. Masalahnya, tebakan selalu berupa satu angka, sedangkan keadaan sawah tidak pernah
sesederhana itu.

Sekarang bagian yang membuat penasaran. Di dua puluh petak itu ia tidak menemukan
rumput sama sekali. Apakah ia boleh pulang dan berkata bahwa sawahnya bersih? Tentu
tidak. Bisa jadi rumputnya memang tidak ada, bisa jadi ia hanya kebetulan memilih
petak-petak yang terdiri dari tanah kosong. Yang dapat ia katakan hanyalah: kalau
rumputnya cukup banyak, dua puluh petak kosong itu hampir mustahil. Dari situ ia bisa
menghitung batas atas yang jujur, misalnya ``paling banyak tiga dari seratus petak''.

Hal yang sama berlaku ketika guru menilai ulangan beberapa murid dan ketika program
komputer diuji pada beberapa contoh.
Di bab ini kita akan tahu mengapa satu angka tidak pernah cukup, mengapa jawaban jujur
selalu berupa selang, dan mengapa selang itu melebar ketika buktinya sedikit.
\end{untukanak}

\section{Penaksir, bias, dan galat baku}
\tingkat{dasar}

Setiap kesimpulan statistika dimulai dari satu keputusan kecil: besaran apa yang kita
hitung dari data. Rata-rata, proporsi, simpangan baku sampel, semuanya adalah fungsi
dari data dan tidak memerlukan pengetahuan tentang parameter yang dicari. Kalimat
sehari-harinya: penaksir adalah juru ukur yang menebak tinggi pohon dari panjang
bayangannya, dan kita perlu tahu apakah juru ukur itu sistematis meleset atau hanya
goyah.

\kamus{penaksir}{estimator}{juru ukur yang menebak tinggi pohon dari panjang
bayangannya, bukan dari pohon itu sendiri}

\begin{definisi}[Statistik dan penaksir]
\index{penaksir}
Misalkan $X_1,\dots,X_n$ sampel acak dari distribusi yang bergantung pada parameter
$\theta$. \emph{Statistik} adalah fungsi dari $X_1,\dots,X_n$ yang tidak memuat
$\theta$. Statistik yang dipakai menebak $\theta$ disebut \emph{penaksir}
(\emph{estimator}) dan ditulis $\hat\theta$; nilai yang diperoleh pada satu sampel
disebut \emph{taksiran} (\emph{estimate}).
\end{definisi}

Penaksir adalah variabel acak, sebab ia fungsi dari variabel acak. Karena itu ia punya
harapan dan variansi, dan kedua besaran itulah yang menentukan mutunya. Harapannya
menjawab pertanyaan ``apakah tebakannya meleset ke satu arah terus'', sedangkan
variansinya menjawab ``seberapa jauh tebakan itu bergoyang dari sampel ke sampel''.

\begin{definisi}[Bias dan penaksir tak bias]
\index{tak bias}
Bias penaksir $\hat\theta$ untuk $\theta$ adalah $\operatorname{bias}(\hat\theta)
=\mathbb{E}[\hat\theta]-\theta$. Penaksir disebut \emph{tak bias} bila biasnya nol,
yakni $\mathbb{E}[\hat\theta]=\theta$ untuk setiap nilai $\theta$ yang mungkin.
\end{definisi}

Rata-rata sampel tak bias untuk $\mu$, dan itu sudah dipakai pada pembuktian hukum
bilangan besar di Bab 19, jadi tidak diulang di sini. Yang lebih menarik adalah
penaksir yang dipakai untuk variansi. Dengan pembagi $n$, penaksir
$S_n^2=\frac1n\sum(X_i-\bar X_n)^2$ bias ke bawah; dengan pembagi $n-1$, penaksir
$S^2=\frac1{n-1}\sum(X_i-\bar X_n)^2$ tak bias. Sebabnya, $\sum(X_i-\bar X_n)^2$
kehilangan satu derajat kebebasan karena $\bar X_n$ sendiri dihitung dari data yang
sama, sehingga harapannya $(n-1)\sigma^2$ dan bukan $n\sigma^2$. Pembagi $n-1$
adalah koreksi untuk kehilangan itu, bukan kerapian belaka.

Kualitas penaksir biasanya diringkas oleh satu angka, yaitu galat kuadrat rata-rata.
Kalimat sehari-harinya: seberapa besar kesalahan yang kita harapkan, dihitung dengan
cara yang menghukum kesalahan besar lebih berat.

\begin{definisi}[Galat kuadrat rata-rata]
\index{galat kuadrat rata-rata}
$\operatorname{MSE}(\hat\theta)=\mathbb{E}[(\hat\theta-\theta)^2]$.
\end{definisi}

\begin{teorema}[Dekomposisi MSE]
Untuk penaksir $\hat\theta$ dengan variansi berhingga,
\[
\operatorname{MSE}(\hat\theta)=\operatorname{bias}(\hat\theta)^2
+\operatorname{Var}(\hat\theta).
\]
\end{teorema}

\begin{proof}[Bukti]
Tulis $\hat\theta-\theta=(\hat\theta-\mathbb{E}\hat\theta)
+(\mathbb{E}\hat\theta-\theta)$. Kuadratkan dan ambil harapannya:
\[
\mathbb{E}[(\hat\theta-\theta)^2]
=\mathbb{E}[(\hat\theta-\mathbb{E}\hat\theta)^2]
+2(\mathbb{E}\hat\theta-\theta)\,\mathbb{E}[\hat\theta-\mathbb{E}\hat\theta]
+(\mathbb{E}\hat\theta-\theta)^2 .
\]
Suku tengah lenyap karena $\mathbb{E}[\hat\theta-\mathbb{E}\hat\theta]=0$, sedangkan
suku pertama adalah variansi dan suku ketiga adalah bias dikuadratkan.
\end{proof}

Identitas itu menjelaskan mengapa ``tak bias'' bukan gelar juara. Penaksir tak bias
boleh bervariansi besar, dan penaksir berbias sedikit boleh bervariansi jauh lebih
kecil sehingga MSE-nya lebih rendah. Untuk menaksir $\mu$ dari $n$ pengamatan,
statistik $X_1$ tak bias dengan variansi $\sigma^2$, sedangkan $\bar X_n$ tak bias
dengan variansi $\sigma^2/n$; pada $n=20$ selisihnya dua puluh kali lipat.

Contoh kedua menunjukkan bias yang harus dihitung, karena ia dipakai setiap hari
untuk proporsi. Bila $K\sim\mathrm{Bin}(n,p)$ dan $\hat p=K/n$, maka
$\hat p(1-\hat p)$ adalah penaksir alami bagi $p(1-p)$, yaitu variansi satu percobaan.
Harapannya
\[
\mathbb{E}[\hat p(1-\hat p)]
=\mathbb{E}[\hat p]-\mathbb{E}[\hat p^2]
=p-\Bigl(p^2+\frac{p(1-p)}{n}\Bigr)
=p(1-p)\Bigl(1-\frac1n\Bigr),
\]
sehingga biasnya $-p(1-p)/n$. Untuk $p=0{,}3$ dan $n=20$ nilainya $p(1-p)=0{,}21$
dan biasnya $-0{,}0105$, yakni lima persen ke bawah. Pada $n=200$ biasnya tinggal
$-0{,}00105$. Bias yang mengecil sebanding $1/n$ ini khas: ia hilang di tak hingga,
tetapi hadir pada setiap sampel berhingga.

Sifat kedua yang dituntut dari penaksir adalah bahwa ia benar-benar menuju sasaran
ketika data bertambah. Kalimat sehari-harinya: dengan cukup banyak bukti, tebakan yang
salah menjadi semakin mustahil, bukan hanya semakin jarang.

\begin{definisi}[Konsisten]
\index{konsisten}
Barisan penaksir $\hat\theta_n$ \emph{konsisten} untuk $\theta$ bila untuk setiap
$\varepsilon>0$ berlaku $P(|\hat\theta_n-\theta|\ge\varepsilon)\to0$ ketika
$n\to\infty$.
\end{definisi}

\begin{proposisi}[MSE menuju nol cukup untuk konsistensi]
\index{konsisten}
Bila $\operatorname{MSE}(\hat\theta_n)\to0$, maka $\hat\theta_n$ konsisten untuk
$\theta$.
\end{proposisi}

\begin{proof}[Bukti]
Menurut ketaksamaan Chebyshev, yang sudah dibuktikan di Bab 19,
\[
P(|\hat\theta_n-\theta|\ge\varepsilon)
\le\frac{\mathbb{E}[(\hat\theta_n-\theta)^2]}{\varepsilon^2}
=\frac{\operatorname{MSE}(\hat\theta_n)}{\varepsilon^2}.
\]
Untuk $\varepsilon$ tetap, ruas kanan menuju nol, dan ketaksamaan itu menutup
pembuktian.
\end{proof}

Perhatikan bahwa Chebyshev memakai MSE, bukan variansi. Suku bias dikuadratkan, jadi
penaksir berbias tetap konsisten asalkan biasnya menuju nol. Sebaliknya penaksir $X_1$
tak bias untuk setiap $n$ tetapi tidak pernah memakai sisa data, sehingga tidak
konsisten.

Besaran berikutnya dipakai terus sampai akhir bab. Variansi penaksir mengatakan
seberapa jauh tebakan bergoyang, dan akarnya punya nama sendiri.

\kamus{galat baku}{standard error}{kisaran goyangan satu tebakan, bukan lebar seluruh
data}

\begin{definisi}[Galat baku]
\index{galat baku}
\emph{Galat baku} penaksir $\hat\theta$ adalah $\operatorname{se}(\hat\theta)
=\sqrt{\operatorname{Var}(\hat\theta)}$. Bila rumusnya memuat parameter yang tidak
diketahui, galat baku \emph{tertaksir} diperoleh dengan mengganti parameter itu oleh
taksirannya.
\end{definisi}

Untuk proporsi, $\operatorname{se}(\hat p)=\sqrt{p(1-p)/n}$, dan versi tertaksirnya
$\sqrt{\hat p(1-\hat p)/n}$. Untuk data $33$ dari $50$, $\hat p=0{,}66$ dan
\[
\sqrt{\frac{0{,}66\cdot0{,}34}{50}}=\sqrt{0{,}004488}=0{,}066993 .
\]
Untuk dua belas pengukuran waktu tunda (dalam milidetik),
\begin{align*}
&18{,}2,\ 19{,}1,\ 17{,}8,\ 20{,}3,\ 18{,}9,\ 19{,}7,\\
&18{,}4,\ 21{,}0,\ 17{,}5,\ 19{,}9,\ 18{,}6,\ 20{,}1,
\end{align*}
rata-ratanya $19{,}125$ dan simpangan baku sampelnya $1{,}082191$. Galat bakunya
$1{,}082191/\sqrt{12}=0{,}312401$.
Angka itu tidak mengatakan bahwa rata-rata sebenarnya berjarak $0{,}312$ dari
$19{,}125$; ia mengatakan bahwa tebakan seperti ini, bila diulang, bergoyang dengan
skala sekitar $0{,}31$. Selang kepercayaan adalah cara mengubah kalimat terakhir
menjadi kalimat yang bisa dipakai.

\section{Interval kepercayaan untuk rata-rata}
\tingkat{menengah}

Satu angka tanpa ukuran goyangan mudah disalahgunakan. Yang dibutuhkan sepasang angka
yang, dengan peluang besar, mengurung parameter. Kalimat sehari-harinya: kita tidak
menembak satu titik, kita memasang jaring dan berjanji jaring itu menangkap ikan yang
kita cari pada sebagian besar lemparan.

\kamus{interval kepercayaan}{confidence interval}{jaring yang dipasang dengan lebar
tertentu, dan janji bahwa jaring itu menangkap sasaran pada sebagian besar lemparan}

\begin{definisi}[Interval kepercayaan dan cakupan]
\index{interval kepercayaan}
Misalkan $C(X_1,\dots,X_n)$ adalah selang yang dihitung dari data. $C$ disebut
\emph{interval kepercayaan} bertaraf $1-\alpha$ untuk $\theta$ bila
$P_\theta(\theta\in C)\ge1-\alpha$ untuk setiap $\theta$. Peluang itu disebut
\emph{cakupan} (\emph{coverage}) prosedurnya.
\end{definisi}

Perhatikan letak peluangnya: $P_\theta$ dihitung sebelum data diambil, dan yang acak
adalah $C$, bukan $\theta$. Ini tempat kesalahan tafsir paling sering terjadi, dan kita
kembali ke sana sesudah rumusnya ada.

\begin{teorema}[Interval-z]
\index{interval kepercayaan}
Misalkan $X_1,\dots,X_n$ bebas dan berdistribusi $N(\mu,\sigma^2)$ dengan $\sigma$
diketahui. Untuk $z=z_{1-\alpha/2}$, selang
\[
C=\Bigl[\bar X_n-z\frac{\sigma}{\sqrt n},\ \bar X_n+z\frac{\sigma}{\sqrt n}\Bigr]
\]
mempunyai cakupan tepat $1-\alpha$. Untuk distribusi lain dengan variansi berhingga,
cakupannya menuju $1-\alpha$ ketika $n\to\infty$.
\end{teorema}

\begin{proof}[Bukti]
Misalkan $Z=\sqrt n(\bar X_n-\mu)/\sigma$. Untuk sampel normal, $Z\sim N(0,1)$ tepat;
untuk distribusi lain, $Z$ menuju $N(0,1)$ menurut teorema limit pusat di Bab 19.
Ketaksamaan $\mu\in C$ setara dengan $|Z|\le z$, karena
\[
|\bar X_n-\mu|\le z\frac{\sigma}{\sqrt n}
\iff \frac{\sqrt n\,|\bar X_n-\mu|}{\sigma}\le z .
\]
Jadi $P(\mu\in C)=P(|Z|\le z)=1-\alpha$ pada kasus normal, dan menuju $1-\alpha$
pada kasus umum.
\end{proof}

Bila $\sigma$ tidak diketahui, seperti pada hampir semua data nyata, penggantian
$\sigma$ oleh $s$ menambah ketidakpastian dan menebalkan ekor. Hasilnya memakai
distribusi $t$ dengan $n-1$ derajat kebebasan, diterbitkan Gosset pada 1908 dengan nama
samaran Student:
\[
C=\Bigl[\bar X_n-t_{n-1,1-\alpha/2}\frac{s}{\sqrt n},\
\bar X_n+t_{n-1,1-\alpha/2}\frac{s}{\sqrt n}\Bigr],
\]
dengan cakupan tepat $1-\alpha$ bila sampelnya normal. Buktinya menuntut distribusi
bersama $\bar X_n$ dan $s$ yang bebas, dan ada di Wasserman, \emph{All of Statistics},
Bab 6. Untuk $n=12$ data waktu tunda tadi, $t_{11,0{,}975}=2{,}200985$, sehingga
setengah lebarnya $2{,}200985\cdot0{,}312401=0{,}687591$ dan selangnya
$(18{,}4374;\ 19{,}8126)$. Dengan $z=1{,}96$ setengah lebarnya $0{,}612296$: mengakui
$\sigma$ tidak diketahui melebarkan selang dua belas persen pada $n=12$, dan satu
persen pada $n=100$.

Sekarang tafsirannya, dan di sini kita harus lambat. Sesudah data diambil,
$C$ bukan lagi variabel acak: ia selang tetap $(18{,}4374;\ 19{,}8126)$. Rata-rata
sebenarnya $\mu$ juga bilangan tetap, hanya saja tidak diketahui. Kalimat
``peluang $\mu$ berada di dalam $(18{,}4374;\ 19{,}8126)$ sama dengan $0{,}95$''
karena itu tidak bermakna dalam kerangka ini: tidak ada peluang yang tersisa, sebab
tidak ada yang acak. Yang berpeluang $0{,}95$ adalah \emph{prosedurnya}, dan itu
pernyataan tentang percobaan yang belum dijalankan.

Gambar~\ref{fig:dua-puluh-selang} menjalankan percobaan itu dua puluh kali, dengan
$\mu=19{,}125$ dan $\sigma=1$ yang diketahui: sembilan belas selang memuat $\mu$, satu
tidak.

\begin{figure}[htbp]
\centering
\begin{tikzpicture}[x=4.15cm,y=0.235cm,font=\scriptsize]
  \draw[->,layergrey] (17.80,0)--(20.22,0);
  \draw[layergrey,densely dashed,thick] (19.125,0)--(19.125,21.4);
  \node[above,font=\scriptsize] at (19.125,21.4) {$\mu=19{,}125$};
  \foreach \y/\a/\b in {
    1/18.5955/19.7271,
    2/18.5211/19.6527,
    3/18.7441/19.8757,
    4/18.5895/19.7211,
    5/18.4046/19.5362,
    6/18.6636/19.7952,
    7/18.9356/20.0672,
    8/18.8326/19.9642,
    9/18.3561/19.4876,
    10/18.1939/19.3255,
    11/18.3793/19.5109,
    12/18.5711/19.7027,
    14/18.4960/19.6276,
    15/18.1995/19.3311,
    16/18.3478/19.4794,
    17/18.4021/19.5337,
    18/18.4679/19.5995,
    19/18.6780/19.8096,
    20/18.8602/19.9917
  } {
    \draw[very thick,layerindigo] (\a,\y)--(\b,\y);
    \draw[layerindigo] (\a,\y-0.22)--(\a,\y+0.22);
    \draw[layerindigo] (\b,\y-0.22)--(\b,\y+0.22);
  }
  \draw[very thick,accent] (17.8880,13)--(19.0196,13);
  \draw[accent] (17.8880,12.7)--(17.8880,13.3);
  \draw[accent] (19.0196,12.7)--(19.0196,13.3);
  \node[accent,anchor=west,font=\scriptsize] at (17.82,0.62) {satu selang meleset};
  \foreach \t in {18,19,20} {\draw[layergrey] (\t,-0.2)--(\t,0.2);}
  \foreach \t in {18,19,20} {\node[below,font=\scriptsize] at (\t,-0.45) {\t};}
\end{tikzpicture}
\caption{Dua puluh selang kepercayaan $95$ persen untuk $\mu$, masing-masing dari satu
sampel berukuran $n=12$ dengan $\sigma=1$ diketahui, sehingga setengah lebarnya
$1{,}96/\sqrt{12}=0{,}5658$ dan lebarnya $1{,}1316$. Garis putus-putus tegak adalah
$\mu=19{,}125$ yang sebenarnya. Sembilan belas selang memuatnya; selang pada baris
ke-13, yaitu $(17{,}8880;\,19{,}0196)$, tidak. Satu kegagalan dalam dua puluh adalah
angka yang wajar untuk janji $95$ persen, bukan tanda prosedur yang rusak.}
\label{fig:dua-puluh-selang}
\end{figure}

\begin{catatan}
Tiga kalimat yang benar, dan satu yang salah. Benar: ``prosedur ini menangkap $\mu$
pada $95$ persen percobaan yang diulang''. Benar: ``kita $95$ persen yakin $\mu$
berada di antara $18{,}44$ dan $19{,}81$'', asalkan ``yakin'' dibaca sebagai ringkasan
percobaan berulang. Benar: ``prosedur ini salah pada sekitar lima dari seratus
kasus''. Salah: ``peluang $\mu$ ada di selang ini $0{,}95$''. Kalimat terakhir
berpura-pura mengetahui sesuatu tentang $\mu$ yang tidak diberikan data.
\end{catatan}

Lebar selang menyusut sebanding $1/\sqrt n$, bukan $1/n$: memangkasnya menjadi separuh
menuntut empat kali data, menjadi sepersepuluh menuntut seratus kali. Aturan ini dapat
dihitung sebelum percobaan dimulai, sebab setengah lebar $z\sigma/\sqrt n\le d$ setara
dengan $n\ge(z\sigma/d)^2$; dengan $\sigma=1$, $z=1{,}96$, dan $d=0{,}1$, dibutuhkan
$n\ge385$.

\begin{tangan}
Siapkan dua puluh kartu, tulis angka $1$ sampai $20$ pada kartu, satu angka per kartu.
Kocok, ambil lima kartu tanpa melihat, jumlahkan, lalu bagi lima. Tulis hasilnya pada
satu baris kertas. Kembalikan kelima kartu, kocok lagi, ulangi. Lakukan tiga puluh kali,
selalu lima kartu, dan susun angka-angka itu sebagai titik pada satu garis bilangan.

Yang harus terlihat: titik-titik itu tidak menumpuk di satu angka. Rata-rata seluruh
kartu sebenarnya $10{,}5$, dan titikmu menyebar di sekitarnya, ada yang $8$, ada yang
$13$. Sekarang ulangi seluruh pekerjaan dengan sepuluh kartu per pengambilan, tiga puluh
kali lagi. Sebaran kedua lebih sempit, tetapi tidak setengahnya: ia menyempit kira-kira
sebanding akar dua, sebab banyak kartu berlipat dua. Tempel kedua baris titik itu
berdampingan, dan lihat sendiri bahwa menggandakan pekerjaan hanya memangkas goyangan
sekitar tiga puluh persen.
\end{tangan}

\section{Proporsi: Wald, Wilson, dan aturan tiga}
\tingkat{menengah}

Untuk proporsi, rumus yang paling banyak dipakai adalah versi paling sederhana dari
teorema sebelumnya: ganti $\bar X_n$ oleh $\hat p$ dan $\sigma/\sqrt n$ oleh galat
baku proporsi. Kalimat sehari-harinya: kita pasang selang simetris di sekitar proporsi
yang teramati, selebar dua kali galat bakunya.

\begin{definisi}[Interval Wald]
\index{interval wald}
Interval Wald bertaraf $1-\alpha$ untuk proporsi adalah
\[
\hat p\pm z\sqrt{\frac{\hat p(1-\hat p)}{n}},\qquad z=z_{1-\alpha/2}.
\]
\end{definisi}

Rumus itu wajar selama $\hat p$ jauh dari nol dan satu. Kesulitannya muncul tepat di
tempat data paling sedikit, dan bentuknya tidak halus. Bila $\hat p=0$, rumus itu
memberi selang $[0;\,0]$ selebar nol: sebuah klaim percaya diri penuh bahwa laju
kejadiannya nol, padahal yang kita punya hanya ketiadaan kejadian pada $n$ percobaan.
Bila $\hat p=1/20$, batas bawahnya $-0{,}0455$, sebuah bilangan yang bukan peluang.

Penyebabnya bukan pembulatan. Wald memakai $\hat p$ untuk dua tugas sekaligus:
sebagai titik tengah dan sebagai penaksir variansi. Ketika $\hat p$ mendekati nol,
variansi taksirannya ikut mengecil, padahal ketidakpastian sebenarnya justru besar pada
skala itu.

Jalan keluarnya adalah membalik pertanyaannya: alih-alih menganggap $\hat p$ menyebar
normal di sekitar $p$ yang benar, kita tanyakan nilai $p$ mana saja yang masih masuk
akal dengan data teramati. Kriterianya, jarak ternormalisasi antara $\hat p$ dan $p$
tidak melampaui $z$, dengan variansi di penyebut memakai $p$, bukan $\hat p$.

\kamus{interval wilson}{Wilson interval}{selang yang dicari dengan bertanya nilai mana
yang masih masuk akal, bukan dengan memasang selang simetris di sekitar tebakan}

\begin{teorema}[Interval Wilson]
\index{interval wilson}
Untuk $0\le k\le n$, $\hat p=k/n$, dan $z>0$, himpunan $p\in[0,1]$ yang memenuhi
\[
n(\hat p-p)^2\le z^2p(1-p)
\]
adalah selang tertutup dengan ujung
\[
\frac{\hat p+\dfrac{z^2}{2n}\ \pm\
z\sqrt{\dfrac{\hat p(1-\hat p)}{n}+\dfrac{z^2}{4n^2}}}{1+\dfrac{z^2}{n}} .
\]
Kedua ujungnya selalu terletak di $[0,1]$, dan selangnya tidak pernah kosong.
\end{teorema}

\begin{proof}[Bukti]
Susun ulang ketaksamaan menjadi $(n+z^2)p^2-(2n\hat p+z^2)p+n\hat p^2\le0$.
Diskriminannya
\[
(2n\hat p+z^2)^2-4(n+z^2)n\hat p^2
=4z^2\Bigl(n\hat p(1-\hat p)+\frac{z^2}{4}\Bigr)\ge0,
\]
jadi akarnya real dan rumus akar persamaan kuadrat memberi bentuk pada pernyataan
teorema setelah pembilang dan penyebut dibagi $2n$.

Untuk letaknya, misalkan $f(p)=n(\hat p-p)^2-z^2p(1-p)$. Koefisien $p^2$ adalah
$n+z^2>0$, jadi $\{f\le0\}$ adalah selang $[r_-,r_+]$. Karena
$f(0)=n\hat p^2\ge0$, titik $0$ tidak berada di bagian dalam selang itu, sehingga
$r_-\ge0$; karena $f(1)=n(1-\hat p)^2\ge0$, titik $1$ juga tidak, sehingga
$r_+\le1$. Terakhir $f(\hat p)=-z^2\hat p(1-\hat p)\le0$ menempatkan $\hat p$ di dalam
selang, jadi selang itu tidak kosong.
\end{proof}

Sekarang bandingkan keduanya pada data yang sama. Tabel berikut memakai $z=1{,}96$ dan
selang kepercayaan $95$ persen; semua angka dihitung dengan rumus di atas dan diperiksa
ulang dengan Python.

\begin{center}
\small
\begin{tabular}{lrrrr}
\hline
data & $\hat p$ & Wald $95\%$ & Wilson $95\%$ & lebar Wilson \\
\hline
$33/50$  & 0,660 & $(0{,}5287;\ 0{,}7913)$ & $(0{,}5215;\ 0{,}7756)$ & 0,254 \\
$53/100$ & 0,530 & $(0{,}4322;\ 0{,}6278)$ & $(0{,}4329;\ 0{,}6249)$ & 0,192 \\
$13/20$  & 0,650 & $(0{,}4410;\ 0{,}8590)$ & $(0{,}4329;\ 0{,}8188)$ & 0,386 \\
$1/20$   & 0,050 & $(-0{,}0455;\ 0{,}1455)$ & $(0{,}0089;\ 0{,}2361)$ & 0,227 \\
$0/20$   & 0,000 & $(0;\ 0)$ & $(0;\ 0{,}1611)$ & 0,161 \\
\hline
\end{tabular}
\end{center}

Dua baris pertama hampir tidak berbeda: ketika $\hat p$ dekat setengah, keduanya sah
dan sepakat sampai dua angka desimal. Baris keempat dan kelima memperlihatkan
pemisahan. Pada $1/20$ selang Wald \emph{lebih sempit} daripada Wilson, $0{,}191$
melawan $0{,}227$, dan kesempitan itulah tandanya salah: batas bawahnya keluar dari
$[0,1]$. Pada $0/20$ Wald memberi lebar nol, sedangkan Wilson menyisakan ruang sampai
$0{,}1611$. Kesimpulannya bukan bahwa Wilson selalu lebih baik, melainkan bahwa lebar
selang tidak dapat dinilai tanpa cakupannya.

Untuk kasus nol kejadian ada hasil khusus yang lebih tua daripada Wilson dan lebih
mudah diingat. Kalimat sehari-harinya: tidak ada kejadian sama sekali bukan bukti
bahwa laju kejadiannya nol, dan batas atasnya kira-kira tiga dibagi banyak percobaan.

\kamus{aturan tiga}{rule of three}{jika tidak ada kejadian dalam sekian percobaan,
laju kejadiannya masih boleh setinggi tiga dibagi banyak percobaan}

\begin{teorema}[Aturan tiga]
\index{aturan tiga}
Misalkan $K\sim\mathrm{Bin}(n,p)$ dan teramati $K=0$. Batas atas satu sisi pada taraf
$95\%$ adalah $p_{\mathrm{hi}}=1-0{,}05^{1/n}$, dan berlaku
\[
p_{\mathrm{hi}}\le\frac{\ln 20}{n}<\frac{3}{n}.
\]
\end{teorema}

\begin{proof}[Bukti]
Peluang mengamati nol kejadian adalah $(1-p)^n$. Menetapkan peluang itu sama dengan
$0{,}05$ memberi $(1-p_{\mathrm{hi}})^n=0{,}05$, yakni
$p_{\mathrm{hi}}=1-0{,}05^{1/n}=1-e^{-c/n}$ dengan $c=\ln 20=2{,}9957\ldots$ Karena
$1-e^{-x}\le x$ untuk $x\ge0$, diperoleh $p_{\mathrm{hi}}\le c/n<3/n$.
\end{proof}

Untuk $n=100$, aturan tiga memberi $3/100=0{,}03$, batas eksaknya $0{,}0295130$, dan
batas atas Wilson dua sisi pada $k=0$ adalah $z^2/(n+z^2)=0{,}0369935$. Ketiga angka
menjawab pertanyaan yang sama dengan tiga janji berbeda; aturan tiga adalah batas satu
sisi yang eksak, Wilson aproksimasi dua sisi. Yang tidak boleh adalah menyebut angka
tanpa menyebut prosedurnya.

\section{Uji hipotesis dan nilai-p}
\tingkat{menengah}

Sampai di sini pertanyaannya berbentuk ``berapa nilainya''. Pertanyaan berikutnya
berbentuk ``apakah nilainya sama dengan angka tertentu''. Dalam bahasa sehari-hari,
kita memasang tuduhan bahwa tidak ada efek, lalu memeriksa apakah data cukup
mengejutkan untuk membatalkan tuduhan itu.

\begin{definisi}[Hipotesis nol, alternatif, dan taraf]
\index{hipotesis nol}
\emph{Hipotesis nol} $H_0$ adalah pernyataan tentang parameter yang ingin diuji;
\emph{hipotesis alternatif} $H_1$ adalah himpunan nilai yang dianggap benar bila $H_0$
ditolak. \emph{Taraf} $\alpha$ adalah peluang menolak $H_0$ ketika $H_0$ sebenarnya
benar. Aturan penolakan ditentukan oleh satu statistik uji $T$ dan satu daerah tolak
$R$: kita menolak $H_0$ bila $T\in R$.
\end{definisi}

Taraf $\alpha$ dipilih sebelum data dilihat, dan di situlah letak seluruh
kredibilitas prosedur: kita bersedia salah menuduh pada sebagian kecil percobaan, dan
kita menyebutkan berapa besar bagian itu. Nilai-p memindahkan gagasan yang sama ke
satu dataset.

\begin{definisi}[Nilai-p]
\index{nilai-p}
Untuk data $x$ dan statistik uji $T$, nilai-p adalah peluang memperoleh hasil yang
sekurang-kurangnya sejauh hasil teramati bila $H_0$ benar,
\[
\text{nilai-}p=\sup_{\theta\in H_0}P_\theta\bigl(T\text{ sejauh }T(x)\bigr).
\]
Kita menolak $H_0$ pada taraf $\alpha$ bila nilai-p $\le\alpha$.
\end{definisi}

Kalimat ``sekurang-kurangnya sejauh'' perlu ditentukan untuk uji dua sisi, dan
ketentuan yang lazim adalah menjumlahkan kedua ekor: untuk statistik $z$ yang
simetris, nilai-p $=2(1-\Phi(|z|))$. Definisi itu membuat nilai-p menjadi variabel
acak, dan sifatnya di bawah $H_0$ dapat dinyatakan tepat.

\begin{teorema}[Nilai-p seragam di bawah $H_0$]
\index{nilai-p}
Bila $H_0$ benar dan $T$ berdistribusi kontinu dengan fungsi sebaran $F_0$, maka
$P_{H_0}(\text{nilai-}p\le\alpha)=\alpha$ untuk setiap $\alpha\in(0,1)$. Akibatnya
$P_{H_0}(\text{nilai-}p\le0{,}05)=0{,}05$: uji dengan aturan ``tolak bila nilai-p
$\le\alpha$'' mempunyai taraf tepat $\alpha$.
\end{teorema}

\begin{proof}[Bukti]
Untuk uji satu sisi kanan, nilai-p $=1-F_0(T)$, sehingga
\[
P_{H_0}(1-F_0(T)\le\alpha)
=P_{H_0}\bigl(F_0(T)\ge1-\alpha\bigr)
=P_{H_0}\bigl(T\ge F_0^{-1}(1-\alpha)\bigr)=\alpha,
\]
karena $F_0(T)$ berdistribusi seragam pada $(0,1)$ bila $F_0$ kontinu. Untuk uji dua
sisi, nilai-p bergantung pada $|T|$ dan pembuktiannya memakai langkah yang sama pada
sebaran $|T|$.
\end{proof}

Keseragaman itu punya akibat yang sering terlewat: di bawah $H_0$ nilai-p tersebar
rata, jadi nilai-p $=0{,}04$ muncul pada empat persen percobaan yang hipotesis nolnya
benar. Yang membedakan dua percobaan bukan nilai-p tunggal, melainkan daya
masing-masing.

\begin{bacakeras}
Dua suara, bergantian.

``Nilai p nol koma nol empat. Berarti peluang tidak ada efek tinggal empat persen.''

``Bukan. Nilai p menghitung data, bukan tuduhan. Seandainya tidak ada efek, dan
percobaan ini diulang seratus kali, sekitar empat kali data akan tampak sejauh ini.''

``Jadi yang dihitung data.''

``Ya, data. Tuduhan hanya tempat berpijak.''

``Dan kalau nilai p besar?''

``Berarti data tidak mengejutkan. Diamnya timbangan bukan bukti bahwa barangnya ringan.''

Ucapkan sekali lagi, pelan: nilai p adalah peluang data, bukan peluang tuduhan. Ulangi
baris itu sampai terasa wajar di mulut: peluang data, bukan peluang tuduhan.
\end{bacakeras}


\begin{resep}
\begin{enumerate}
  \item Tulis $H_0$ dan $H_1$ sebagai pernyataan tentang satu parameter, sebelum
        menghitung apa pun.
  \item Tetapkan $\alpha$ dan daerah tolaknya, lalu berhenti sejenak. Jangan melihat
        data dulu.
  \item Hitung statistik uji dengan galat baku di bawah $H_0$, lalu nilai-p.
  \item Bandingkan nilai-p dengan $\alpha$ yang sudah ditetapkan, dan laporkan
        selangnya bersama keputusan itu.
\end{enumerate}
Peringatan urutan: langkah 2 harus selesai sebelum langkah 3. Bila $\alpha$ dipilih
sesudah data terlihat, $\alpha$ bukan lagi janji melainkan catatan hasil, dan langkah 4
kehilangan artinya. Yang tidak boleh ditukar: langkah 2 dan langkah 3.
\end{resep}

Sekarang contoh dengan angka. Untuk data $33$ dari $50$ dan $H_0\colon p=0{,}5$,
statistik $z$ memakai galat baku di bawah $H_0$, yaitu $\sqrt{0{,}25/50}=0{,}070711$:
\[
z=\frac{0{,}66-0{,}5}{0{,}070711}=2{,}2627,\qquad
\text{nilai-}p=2(1-\Phi(2{,}2627))=0{,}0237 .
\]
Uji binomial eksak atas aturan yang sama, $P(X\ge33)+P(X\le17)$ untuk
$X\sim\mathrm{Bin}(50;0{,}5)$, memberi $0{,}0328$. Kedua angka menolak $H_0$ pada
taraf lima persen, tetapi aproksimasi normal melaporkan nilai-p hampir tiga puluh
persen lebih kecil. Pada $n=50$ selisih sebesar itu cukup membalik keputusan bila
tarafnya $0{,}03$; laporan yang baik menyebutkan uji mana yang dipakai, bukan hanya
nilai-p-nya.

Gambar~\ref{fig:dua-ekor} menaruh kedua batas itu pada sebaran $z$ di bawah $H_0$: dua
ekor seluas $0{,}0250$ masing-masing, dan statistik teramati $2{,}2627$ di dalam ekor
kanan.

\begin{figure}[htbp]
\centering
\begin{tikzpicture}
\begin{axis}[
  width=11.2cm, height=5.4cm,
  xmin=-4, xmax=4, ymin=0, ymax=0.46,
  axis lines=left,
  xtick={-1.96,0,1.96},
  xticklabels={$-1{,}96$,$0$,$1{,}96$},
  ytick=\empty,
  tick label style={font=\small},
  xlabel={nilai $z$ bila $H_0$ benar},
  xlabel style={font=\small},
]
\addplot[fill=accent!25,draw=none,domain=-3.95:-1.96,samples=100]
  {0.39894228*exp(-0.5*x^2)} \closedcycle;
\addplot[fill=accent!25,draw=none,domain=1.96:3.95,samples=100]
  {0.39894228*exp(-0.5*x^2)} \closedcycle;
\addplot[very thick,layerindigo,domain=-3.95:3.95,samples=200]
  {0.39894228*exp(-0.5*x^2)};
\draw[densely dashed,layergrey] (axis cs:-1.96,0)--(axis cs:-1.96,0.0584);
\draw[densely dashed,layergrey] (axis cs:1.96,0)--(axis cs:1.96,0.0584);
\draw[very thick,accent] (axis cs:2.2627,0)--(axis cs:2.2627,0.0308);
\node[font=\scriptsize] (tl) at (axis cs:-2.55,0.135) {$0{,}025$};
\draw[-{Stealth[length=4pt]},layergrey] (tl)--(axis cs:-2.55,0.012);
\node[accent,font=\scriptsize] (tz) at (axis cs:2.05,0.215) {$z=2{,}2627$};
\draw[-{Stealth[length=4pt]},accent] (tz)--(axis cs:2.2627,0.055);
\node[font=\scriptsize] (tr) at (axis cs:3.05,0.135) {$0{,}025$};
\draw[-{Stealth[length=4pt]},layergrey] (tr)--(axis cs:3.05,0.004);
\end{axis}
\end{tikzpicture}
\caption{Sebaran statistik $z=\sqrt n(\hat p-p_0)/\sqrt{p_0(1-p_0)}$ bila $H_0$ benar,
untuk data $33$ dari $50$ dengan $H_0\colon p=0{,}5$. Dua garis putus-putus di
$\pm1{,}96$ memotong dua ekor yang diwarnai, masing-masing seluas
$1-\Phi(1{,}96)=0{,}0250$, sehingga taraf dua sisi $0{,}05$. Garis tebal di
$2{,}2627$ adalah statistik yang teramati pada data itu; ekor kanan dari titik itu
seluas $0{,}0118$, dan dua kali luas itu adalah nilai-p $0{,}0237$ yang dilaporkan uji
normal.}
\label{fig:dua-ekor}
\end{figure}

Interval kepercayaan dan uji hipotesis bukan dua prosedur berbeda; keduanya dua cara
membaca himpunan yang sama.

\begin{proposisi}[Kesetaraan selang dan uji]
\index{interval kepercayaan}
Untuk selang Wald dan uji-$z$ dua sisi yang memakai galat baku $\hat p$, berlaku
$p_0$ termuat dalam selang Wald bertaraf $1-\alpha$ jika dan hanya jika $H_0\colon
p=p_0$ tidak ditolak pada taraf $\alpha$. Secara umum, selang Wilson adalah himpunan
$p_0$ yang tidak ditolak oleh uji skor pada taraf $\alpha$.
\end{proposisi}

\begin{proof}[Bukti]
Untuk pasangan Wald, kedua pernyataan memakai ketaksamaan yang sama, yaitu
$|\hat p-p_0|\le z\sqrt{\hat p(1-\hat p)/n}$: yang pertama membacanya sebagai
``$p_0$ berada di dalam selang'', yang kedua sebagai ``statistik uji tidak melewati
$z$''. Untuk Wilson, himpunan $\{p_0:n(\hat p-p_0)^2\le z^2p_0(1-p_0)\}$ adalah
definisi uji skor yang dibalik, dan teorema Wilson sudah menghitung ujungnya.
\end{proof}

\section{Galat tipe I, galat tipe II, dan daya}
\tingkat{lanjut}

Setiap keputusan biner dapat salah dengan dua cara, dan keduanya tidak simetris.
Kalimat sehari-harinya: alarm kebakaran dapat berbunyi padahal tidak ada api, dan
dapat diam padahal ada api. Kesalahan pertama membuat kita keluar rumah tanpa sebab;
kesalahan kedua jauh lebih mahal.

\begin{definisi}[Galat tipe I, galat tipe II, dan daya]
\index{galat tipe i}
\emph{Galat tipe I} adalah menolak $H_0$ ketika $H_0$ benar; peluangnya adalah taraf
$\alpha$. \emph{Galat tipe II} adalah tidak menolak $H_0$ ketika $H_1$ benar;
peluangnya ditulis $\beta$. \emph{Daya} uji adalah $1-\beta$, yaitu peluang menolak
$H_0$ ketika efek yang dicari benar-benar ada.
\end{definisi}

Gambar~\ref{fig:dua-galat} menyusun keempat kemungkinan itu sebagai kotak keputusan,
dengan peluang dari hitungan eksak pada $n=20$, $p_0=0{,}5$, dan $p_1=0{,}75$ yang
muncul di Lab komputasi: $\alpha=0{,}0414$, daya $0{,}6172$, jadi $\beta=0{,}3828$.

\begin{figure}[htbp]
\centering
\begin{tikzpicture}[font=\small,
  keputusan/.style={draw=layergrey,thick,rounded corners=2pt,
    minimum width=4.3cm,minimum height=1.6cm,align=center,inner sep=4pt}]
\node[keputusan,fill=layerjade!10] (a) at (0,0)
  {tidak menolak $H_0$\\[3pt] benar, peluang $1-\alpha=0{,}9586$};
\node[keputusan,fill=accent!12] (b) at (5.1,0)
  {menolak $H_0$\\[3pt] galat tipe I, peluang $\alpha=0{,}0414$};
\node[keputusan,fill=accent!12] (c) at (0,-2.4)
  {tidak menolak $H_0$\\[3pt] galat tipe II, peluang $\beta=0{,}3828$};
\node[keputusan,fill=layerjade!10] (d) at (5.1,-2.4)
  {menolak $H_0$\\[3pt] benar, daya $1-\beta=0{,}6172$};
\node[font=\scriptsize,above=2pt] at (0,1.0) {$H_0$ benar: $p=0{,}5$};
\node[font=\scriptsize,above=2pt] at (5.1,1.0) {$H_1$ benar: $p=0{,}75$};
\end{tikzpicture}
\caption{Dua jenis kesalahan sebagai empat kotak keputusan. Baris atas: keadaan
sebenarnya $H_0$ benar; baris bawah: $H_1$ benar dengan $p=0{,}75$. Kolom kiri: uji
tidak menolak; kolom kanan: uji menolak. Kotak berwarna merah adalah dua kesalahan yang
tidak dapat ditekan bersamaan: mengecilkan $\alpha=0{,}0414$ akan memperbesar
$\beta=0{,}3828$, dan sebaliknya. Angka dihitung dari sebaran binomial eksak pada
$n=20$ dengan aturan tolak $k\le5$ atau $k\ge15$.}
\label{fig:dua-galat}
\end{figure}

Daya bukan sifat uji saja, melainkan sifat uji bersama ukuran sampel dan besar efek.
Karena itu ia harus dihitung untuk satu nilai $p_1$ yang dianggap penting, bukan
secara umum. Untuk uji dua sisi dengan $H_0\colon p=0{,}5$, $n=100$, dan aturan tolak
bila $|z|>1{,}96$, daerah tolaknya adalah $k\le40$ atau $k\ge60$. Hitungan eksak atas
binomial memberi ukuran $0{,}0569$, bukan $0{,}05$, dan daya $0{,}8750$ pada
$p_1=0{,}65$ sedangkan aproksimasi normal memberi $0{,}8508$. Aproksimasi normal,
yakni
\[
1-\beta\approx\Phi\Bigl(\frac{(p_1-p_0)\sqrt n}{\sqrt{p_0(1-p_0)}}-z\Bigr),
\]
cukup dekat untuk perencanaan, tetapi selisih dua setengah poin itu nyata dan berasal
dari kediskretan data.

Pelajaran pertama: taraf nominal jarang tercapai tepat. Pada $n=20$ aturan yang sama
menolak bila $k\le5$ atau $k\ge15$, dengan ukuran eksak $0{,}0414$; pada $n=100$
ukurannya $0{,}0569$. Uji diskret tidak dapat memilih ambang di antara dua bilangan
bulat, jadi taraf sebenarnya menggantung pada $n$ dan $p_0$. Pelajaran kedua: pada
$n=20$ daya terhadap $p_1=0{,}75$ hanya $0{,}6172$, sehingga hampir empat dari sepuluh
percobaan berukuran dua puluh gagal menemukannya.

Angka terakhir itu membuka salah tafsir yang paling mahal: membaca ``nilai-p
$>0{,}05$'' sebagai ``tidak ada efek''. Pertama, nilai-p adalah
$P(\text{data}\mid H_0)$, bukan $P(H_0\mid\text{data})$; keduanya peluang bersyarat
dengan syarat tertukar, dan pertukaran itu sudah kita temui pada uji diagnostik di
Bab 17. Kedua, nilai-p besar adalah pernyataan tentang data, bukan tentang dunia: data
yang sedikit tidak dapat mengejutkan siapa pun. Ketiga, kegagalan menolak dapat terjadi
justru karena ujinya lemah, seperti pada $n=20$ tadi.

Contoh kecil membuat ketiganya konkret. Sepuluh keberhasilan dari dua puluh percobaan
terhadap $H_0\colon p=0{,}5$ memberi nilai-p dua sisi tepat $1$: data itu berada
persis di tengah. Apakah itu bukti bahwa $p=0{,}5$? Tidak. Selang Wilson $95$ persen
untuk data yang sama adalah $(0{,}2993;\ 0{,}7007)$. Jadi laju sebenarnya boleh
$0{,}30$, boleh $0{,}70$, dan data ini tidak dapat membedakannya dari $0{,}50$.
Nilai-p $=1$ hanya berarti ``tidak ada arah penyimpangan'', dan itu jauh lebih sedikit
daripada yang biasa dibaca orang.

\begin{tantangan}
Dua versi pendeteksi diuji pada \mbox{$100\,000$} contoh masing-masing. Versi A menangkap
\mbox{$66\,000$}, versi B menangkap \mbox{$65\,000$}. Uji dua proporsi memberi nilai-p sekitar
$0{,}0000026$, dan tim menulis di laporan: ``A lebih baik, $p<0{,}001$''. Apakah
kalimat itu didukung data? Jawab dengan satu selang, bukan dengan nilai-p.

Petunjuk 1: nilai-p mengukur kejutan, bukan besar efek. Hitung selang $95$ persen untuk
selisih kedua proporsi, dan perhatikan satuannya.

Petunjuk 2: galat baku selisihnya $0{,}002126$, jadi selangnya terbentang dari
$0{,}0058$ sampai $0{,}0142$, yaitu nol koma enam sampai satu koma empat poin persen.
Bandingkan dengan ambang praktis yang tim itu sendiri pakai: tiga poin persen.

Kalau mentok, kembali ke bagian Galat tipe I, galat tipe II, dan daya.
\end{tantangan}

\begin{untukanak}
Seorang wasit garis bertugas mengangkat bendera bila bola keluar lapangan. Ia hanya
mengangkat bendera kalau melihat sendiri bola itu keluar, dengan cukup jelas. Bila ia
tidak mengangkat bendera, apakah itu berarti bola selalu di dalam? Tentu tidak. Bisa
jadi bola keluar, tetapi pandangannya terhalang pemain lain atau bolanya sangat cepat.

Cara ia bekerja mirip dengan cara sebuah uji hipotesis bekerja. Ia diberi satu aturan
yang harus diputuskan sebelum pertandingan: seberapa yakin ia harus sebelum
mengangkat bendera. Kalau aturannya terlalu longgar, ia sering mengangkat bendera
padahal bola di dalam, dan penonton marah. Kalau aturannya terlalu ketat, ia hampir
tidak pernah mengangkat bendera, dan pelanggaran besar lolos begitu saja. Kedua
kesalahan itu berbeda, dan ia tidak bisa menekan keduanya sekaligus hanya dengan
kehati-hatian.

Di bab ini kita akan tahu mengapa ``tidak ada bukti'' bukan sama dengan ``tidak ada
pelanggaran'', mengapa jumlah pengamatan berperan sebesar aturan yang dipilih, dan
mengapa wasit yang tidak pernah mengangkat bendera bukan wasit yang lebih teliti.
\end{untukanak}

\section{Contoh hitung: 33 dari 50 dan 53 dari 100}
\tingkat{menengah}

Kita kembali ke dua pendeteksi pada pembuka bab dan menghitung kedua selangnya sampai
ujung. Semua angka di bawah memakai $z=1{,}96$, sehingga $z^2=3{,}8416$.

Untuk A, $\hat p=0{,}66$ dan $n=50$. Suku pertama di dalam akar Wilson adalah
$0{,}66\cdot0{,}34/50=0{,}2244/50=0{,}004488$; suku kedua
$z^2/(4n^2)=3{,}8416/10000=0{,}00038416$; jumlahnya $0{,}00487216$ dan akarnya
$0{,}069801$. Titik tengahnya
\[
\frac{0{,}66+3{,}8416/100}{1+3{,}8416/50}
=\frac{0{,}698416}{1{,}076832}=0{,}648584,
\]
dan setengah lebarnya $1{,}96\cdot0{,}069801/1{,}076832=0{,}127048$. Jadi selang
Wilson untuk A adalah $(0{,}5215;\ 0{,}7756)$, selebar $0{,}2541$. Selang Wald untuk
data yang sama memakai galat baku $0{,}066993$ dan setengah lebar
$1{,}96\cdot0{,}066993=0{,}131303$, sehingga $(0{,}5287;\ 0{,}7913)$ selebar
$0{,}2626$.

Untuk B, $\hat p=0{,}53$ dan $n=100$. Suku pertama
$0{,}53\cdot0{,}47/100=0{,}002491$, suku kedua $3{,}8416/40000=0{,}00009604$,
jumlahnya $0{,}00258704$ dan akarnya $0{,}050863$. Titik tengahnya
$(0{,}53+0{,}019208)/1{,}038416=0{,}528890$, setengah lebarnya
$1{,}96\cdot0{,}050863/1{,}038416=0{,}096003$, jadi $(0{,}4329;\ 0{,}6249)$ selebar
$0{,}1920$. Wald memberi $(0{,}4322;\ 0{,}6278)$ selebar $0{,}1956$.

Sekarang bandingkan. Selang A dan B bertumpuk pada $(0{,}5215;\ 0{,}6249)$, seluas
$0{,}1034$: sekitar $41$ persen panjang A dan $54$ persen panjang B. Uji dua proporsi
berbasis normal memberi $z=1{,}5175$ dan nilai-p dua sisi $0{,}1291$. Meskipun angka
mentahnya berbeda tiga belas poin, data ini tidak cukup menyimpulkan bahwa A lebih
baik; selang A juga lebih lebar, sebab dibangun dari lima puluh contoh.

\kertas{Hitung ulang bagian A tanpa melihat teks: $0{,}66\cdot0{,}34$, bagi $50$,
tambahkan $3{,}8416/10000$, tarik akarnya, lalu bagi dengan $1{,}076832$ sesudah
dikali $1{,}96$. Cocokkan bahwa hasilnya $0{,}127048$. Lalu jawab: berapa banyak
contoh yang diperlukan agar lebar selang A turun menjadi $0{,}15$, bila $\hat p$
tetap $0{,}66$? Petunjuk: lebar sebanding $1/\sqrt n$, dan Anda sudah tahu lebarnya
pada $n=50$.}

\section{Lab komputasi: cakupan, galat tipe I, dan daya}
\tingkat{lanjut}

Seluruh klaim di atas dapat diperiksa dengan komputer, dan ada satu klaim yang hanya
dapat diperiksa dengan cara itu: cakupan sebenarnya dari sebuah prosedur. Kuncinya
adalah mengetahui $p$ yang benar, sesuatu yang tidak pernah terjadi pada data nyata.
Karena itu kita bekerja dengan simulasi dan dengan hitungan eksak atas binomial.

\begin{lab}
\begin{verbatim}
import numpy as np
from math import sqrt
from scipy.stats import norm, binom

z = norm.ppf(0.975)


def wald(k, n):
    p = k / n
    se = sqrt(p * (1 - p) / n)
    return p - z * se, p + z * se


def wilson(k, n):
    p = k / n
    z2 = z * z
    mid = (p + z2 / (2 * n)) / (1 + z2 / n)
    half = z * sqrt(p * (1 - p) / n + z2 / (4 * n * n)) / (1 + z2 / n)
    return mid - half, mid + half


# 1. lebar interval 95% untuk beberapa proporsi
print("== 1. lebar interval 95% ==")
print("   n    k   phat   lebar Wald   lebar Wilson")
for n in (20, 100):
    for p in (0.00, 0.05, 0.25, 0.50, 0.66):
        k = round(p * n)
        a = wald(k, n)
        b = wilson(k, n)
        print("%4d %4d  %5.2f   %9.4f    %9.4f"
              % (n, k, k / n, a[1] - a[0], b[1] - b[0]))

# 2. cakupan eksak pada p = 0,05
print("== 2. cakupan eksak, p=0.05 ==")
for n in (20, 50, 100):
    for nama, f in (("Wald", wald), ("Wilson", wilson)):
        c = 0.0
        for k in range(n + 1):
            lo, hi = f(k, n)
            if lo <= 0.05 <= hi:
                c += binom.pmf(k, n, 0.05)
        print("n=%3d %-6s cakupan=%.4f" % (n, nama, c))

# 3. cakupan simulasi dengan seed tetap
rng = np.random.default_rng(20261024)
R = 20000
print("== 3. cakupan simulasi, R=%d ==" % R)
for n in (20, 100):
    X = rng.binomial(n, 0.05, R)
    for nama, f in (("Wald", wald), ("Wilson", wilson)):
        ok = 0
        w = 0.0
        for k in X:
            lo, hi = f(int(k), n)
            ok += (lo <= 0.05 <= hi)
            w += hi - lo
        print("n=%3d %-6s cakupan=%.4f  lebar=%.4f" % (n, nama, ok / R, w / R))

# 4. galat tipe I dan daya untuk H0: p=0.5
print("== 4. galat tipe I dan daya, H0 p=0.5 ==")


def tolak(k, n, p0=0.5):
    return abs((k / n - p0) / sqrt(p0 * (1 - p0) / n)) > z


for n in (20, 100):
    rej = [k for k in range(n + 1) if tolak(k, n)]
    lo = max(k for k in rej if k < n / 2)
    hi = min(k for k in rej if k > n / 2)
    print("n=%3d tolak bila k<=%d atau k>=%d" % (n, lo, hi))
    print("      ukuran eksak di p=0.5 : %.4f"
          % sum(binom.pmf(k, n, 0.5) for k in rej))
    for p1 in (0.65, 0.75):
        ex = sum(binom.pmf(k, n, p1) for k in rej)
        zc = (p1 - 0.5) / sqrt(0.25 / n)
        ap = norm.cdf(zc - z) + norm.cdf(-zc - z)
        print("      daya di p=%.2f: eksak %.4f  aproksimasi %.4f"
              % (p1, ex, ap))

# 4b. simulasi tingkat galat tipe I ketika H0 benar
print("== 4b. simulasi galat tipe I, R=%d ==" % R)
for n in (20, 100):
    X0 = rng.binomial(n, 0.5, R)
    tolak0 = np.abs((X0 / n - 0.5) / sqrt(0.25 / n)) > z
    print("n=%3d p=0.5: tolak=%.4f" % (n, tolak0.mean()))

# 5. interval-t untuk 12 pengukuran
x = np.array([18.2, 19.1, 17.8, 20.3, 18.9, 19.7,
              18.4, 21.0, 17.5, 19.9, 18.6, 20.1])
n = len(x)
m = x.mean()
s = x.std(ddof=1)
from scipy.stats import t
tq = t.ppf(0.975, n - 1)
h = tq * s / sqrt(n)
print("== 5. interval-t, n=12 ==")
print("rata-rata=%.6f  s=%.6f  galat baku=%.6f" % (m, s, s / sqrt(n)))
print("t=%.6f  setengah lebar=%.6f  IK=(%.6f, %.6f)" % (tq, h, m - h, m + h))
print("setengah lebar dengan z=1.96: %.6f" % (z * s / sqrt(n)))
\end{verbatim}
\end{lab}

Keluaran yang sebenarnya, apa adanya:

\begin{verbatim}
== 1. lebar interval 95% ==
   n    k   phat   lebar Wald   lebar Wilson
  20    0   0.00      0.0000       0.1611
  20    1   0.05      0.1910       0.2272
  20    5   0.25      0.3795       0.3568
  20   10   0.50      0.4383       0.4014
  20   13   0.65      0.4181       0.3860
 100    0   0.00      0.0000       0.0370
 100    5   0.05      0.0854       0.0902
 100   25   0.25      0.1697       0.1676
 100   50   0.50      0.1960       0.1923
 100   66   0.66      0.1857       0.1826
== 2. cakupan eksak, p=0.05 ==
n= 20 Wald   cakupan=0.6389
n= 20 Wilson cakupan=0.9245
n= 50 Wald   cakupan=0.9199
n= 50 Wilson cakupan=0.9622
n=100 Wald   cakupan=0.8775
n=100 Wilson cakupan=0.9659
== 3. cakupan simulasi, R=20000 ==
n= 20 Wald   cakupan=0.6430  lebar=0.1472
n= 20 Wilson cakupan=0.9224  lebar=0.2193
n=100 Wald   cakupan=0.8783  lebar=0.0829
n=100 Wilson cakupan=0.9667  lebar=0.0884
== 4. galat tipe I dan daya, H0 p=0.5 ==
n= 20 tolak bila k<=5 atau k>=15
      ukuran eksak di p=0.5 : 0.0414
      daya di p=0.65: eksak 0.2457  aproksimasi 0.2687
      daya di p=0.75: eksak 0.6172  aproksimasi 0.6088
n=100 tolak bila k<=40 atau k>=60
      ukuran eksak di p=0.5 : 0.0569
      daya di p=0.65: eksak 0.8750  aproksimasi 0.8508
      daya di p=0.75: eksak 0.9997  aproksimasi 0.9988
== 4b. simulasi galat tipe I, R=20000 ==
n= 20 p=0.5: tolak=0.0430
n=100 p=0.5: tolak=0.0554
== 5. interval-t, n=12 ==
rata-rata=19.125000  s=1.082191  galat baku=0.312401
t=2.200985  setengah lebar=0.687591  IK=(18.437409, 19.812591)
setengah lebar dengan z=1.96: 0.612296
\end{verbatim}

Empat hal perlu dibaca dari keluaran itu. Pertama, bagian 2 dan 3 hampir sama sampai
tiga angka, dan kesamaan itu sendiri adalah pemeriksaan, sebab simulasi dengan
$R=20000$ ulangan hanya dapat mendekati hitungan eksak sampai sekitar
$\sqrt{0{,}25/20000}=0{,}0035$. Pada $p=0{,}05$ dan $n=20$ selisihnya $0{,}0041$ untuk
Wald dan $0{,}0021$ untuk Wilson. Simulasi yang jujur harus melaporkan $R$.

Kedua, cakupan Wald pada $n=20$ dan $p=0{,}05$ adalah $0{,}6389$, bukan $0{,}95$.
Penyebabnya dapat dihitung dengan tangan: peluang mengamati nol kejadian $0{,}3585$,
dan pada $k=0$ selang Wald adalah $[0;0]$, yang tidak memuat $0{,}05$; satu titik itu
saja menghabiskan hampir tiga puluh enam poin cakupan. Wilson hanya gagal bila
$k\ge3$, peluangnya $0{,}0755$, dan kegagalannya tidak total karena selangnya tetap
lebar. Itulah sebabnya ia mencapai $0{,}9245$. Perhatikan juga kolom lebar: rata-rata
Wald $0{,}1472$ lebih kecil daripada Wilson $0{,}2193$, dan justru di situ intinya.
Selang yang lebih sempit bukan selang yang lebih baik bila ia sempit karena salah
tempat.

Ketiga, bagian 4 memisahkan ukuran uji dari taraf nominal: $0{,}0414$ pada $n=20$ dan
$0{,}0569$ pada $n=100$, keduanya untuk aturan $|z|>1{,}96$; bagian 4b mengulangnya
dengan simulasi dan mendarat dekat, $0{,}0430$ dan $0{,}0554$. Daya bergerak jauh lebih
drastis: $0{,}2457$ pada $n=20$ dan $0{,}8750$ pada $n=100$ untuk $p_1=0{,}65$.
Aproksimasi normal tidak selalu di sisi yang aman: pada $n=20$ ia melaporkan $0{,}2687$
sedangkan kebenarannya $0{,}2457$.

Keempat, bagian 5 hanya mengulang interval-$t$ yang sudah dihitung di teks, dengan
selisih setengah lebar $t$ dan $z$ sebesar $0{,}075295$.

Program tidak membuktikan satu pun dari semua itu. Ia memeriksa klaim yang sudah
dibuktikan, dan pada dua tempat ia menunjukkan batas dirinya sendiri: cakupan eksak
memakai distribusi yang kita andaikan benar, sehingga bila modelnya salah, angka
$0{,}9245$ itu tidak berarti apa-apa; dan simulasi hanya mengukur seberapa sering
prosedurnya berhasil pada model itu, bukan seberapa benar modelnya.

\begin{untukdosen}
Urutan mengajar yang saya sarankan: mulai dari rapat pagi pada pembuka bab, bukan dari
definisi penaksir. Tulis $33/50$ dan $53/100$ di papan, minta kelas memilih, lalu
tanyakan apa yang terjadi bila lima puluh contoh itu diganti. Dari kegelisahan itu
definisi penaksir, bias, dan galat baku muncul sebagai jawaban. Dekomposisi MSE
dibuktikan di papan dengan satu baris aljabar. Tafsiran selang kepercayaan masuk
sesudahnya, sebagai bahan kuliah tersendiri, bukan catatan kaki.

Dua miskonsepsi yang paling sering muncul. Pertama, mahasiswa membaca ``$95$ persen''
sebagai peluang bersyarat pada data yang sudah ada, lalu menulis ``peluang $\mu$ ada di
selang ini $0{,}95$''; gejalanya tampak ketika mereka tidak dapat menjelaskan apa yang
acak pada selang yang sudah dihitung. Kedua, mereka menganggap nilai-p besar sebagai
bukti $H_0$ benar, dan akan menghitung nilai-p untuk $10/20$ lalu menyimpulkan bahwa
$p=0{,}5$ terbukti.

Aktivitas papan tulis yang berhasil: bagi kelas menjadi dua kelompok, satu menghitung
selang Wald dan satu Wilson untuk $1/20$, lalu suruh keduanya mempertahankan hasilnya
di depan kelas. Kelompok Wald akan menemukan batas bawah negatif dan harus
menjelaskannya. Tanda bahwa kelas siap lanjut: mereka dapat menyebut satu contoh
ukuran sampel yang membuat nilai-p besar padahal efeknya besar, dan dapat menjelaskan
mengapa hal itu terjadi.
\end{untukdosen}

\section{Ringkasan}
\tingkat{dasar}

\begin{ringkasan}
\begin{itemize}
  \item Penaksir adalah statistik yang menebak parameter, jadi ia variabel acak dengan
        harapan dan variansi. Mutunya diukur oleh
        $\operatorname{MSE}=\operatorname{bias}^2+\operatorname{Var}$; tak bias saja
        tidak menjamin apa pun, dan MSE menuju nol menjamin konsistensi lewat
        Chebyshev.
  \item Galat baku adalah akar variansi penaksir, bukan lebar data. Untuk proporsi
        nilainya $\sqrt{\hat p(1-\hat p)/n}$, mengecil sebanding $1/\sqrt n$,
        sehingga empat kali data hanya memangkasnya menjadi separuh.
  \item Interval kepercayaan $95$ persen adalah pernyataan tentang prosedur, bukan
        tentang parameter: prosedur itu menangkap nilai sebenarnya pada $95$ persen
        percobaan yang diulang. Sesudah data diambil, selang itu tetap dan tidak lagi
        berpeluang.
  \item Interval Wald runtuh di dekat $0$ dan $1$: pada $0/20$ lebarnya nol, dan pada
        $p=0{,}05$ dengan $n=20$ cakupannya hanya $0{,}6389$. Interval Wilson, hasil
        membalik uji skor, mencapai $0{,}9245$ pada setelan yang sama.
  \item Bila tidak ada kejadian dalam $n$ percobaan, batas atas $95$ persen kira-kira
        $3/n$; untuk $n=100$ batas eksaknya $0{,}0295130$ dan Wilson dua sisi
        $0{,}0369935$.
  \item Nilai-p adalah peluang data sejauh itu bila $H_0$ benar, dan di bawah $H_0$
        ia tersebar seragam. Nilai-p besar bukan bukti $H_0$ benar: pada $n=20$ daya
        terhadap $p=0{,}75$ hanya $0{,}6172$.
\end{itemize}
\end{ringkasan}

\section{Latihan}
\tingkat{menengah}

\begin{latihan}
\begin{enumerate}
  \item Nilai klaim ini: ``Karena nilai-p sama dengan $0{,}42$, hipotesis nol
        $p=0{,}5$ terbukti benar.'' Benar, salah, atau belum cukup bukti? Sebutkan
        bagian yang benar, bagian yang keliru, dan satu perhitungan tambahan yang
        akan memperkuat jawaban Anda.
  \item Buktikan bahwa $\mathbb{E}[\hat p(1-\hat p)]=p(1-p)(n-1)/n$ untuk
        $\hat p=K/n$ dengan $K\sim\mathrm{Bin}(n,p)$, lalu tentukan penaksir tak bias
        bagi $p(1-p)$ yang diperoleh dari situ. Periksa hasilnya pada $p=0{,}3$ dan
        $n=20$.
  \item Buktikan bahwa bila $\operatorname{bias}(\hat\theta_n)\to0$ dan
        $\operatorname{Var}(\hat\theta_n)\to0$, maka $\hat\theta_n$ konsisten untuk
        $\theta$. Jelaskan mengapa syarat itu lebih lemah daripada
        $\operatorname{MSE}\to0$ tetapi cukup untuk kesimpulan yang sama.
  \item Hitung selang Wilson dan Wald $95$ persen untuk $7/40$, lalu bandingkan
        lebarnya dengan $13/20$. Jelaskan mengapa lebar yang lebih kecil pada $7/40$
        bukan tanda ketelitian yang lebih tinggi, dan periksa kedua selang dengan
        Python.
  \item Untuk $n=30$ dan $H_0\colon p=0{,}5$ dengan uji dua sisi berbasis normal,
        tentukan daerah tolaknya, hitung ukuran eksaknya, lalu hitung dayanya pada
        $p_1=0{,}7$. Bandingkan daya itu dengan $n=100$ pada nilai $p_1$ yang sama.
  \item Buktikan bahwa selang Wald dan uji-$z$ dua sisi yang memakai galat baku
        $\hat p$ memang setara, dengan menuliskan kedua ketaksamaannya berdampingan.
        Jelaskan mengapa kesetaraan itu tidak berlaku bila uji memakai galat baku di
        bawah $H_0$.
  \item Sebuah laporan menulis: ``tidak ada perbedaan yang signifikan, $p=0{,}31$,
        jadi kedua metode sama baik''. Tulis ulang kalimat itu sehingga jujur, dengan
        menyebut selang, taraf, dan daya. Anda boleh mengarang angkanya asalkan
        konsisten.
  \item Untuk $n=100$ tanpa satu pun kejadian, hitung aturan tiga, batas eksak satu
        sisi, dan batas atas Wilson. Jelaskan mengapa ketiga angka berbeda, lalu
        sebutkan mana yang boleh dipakai untuk kalimat ``paling tinggi $5$ persen''.
  \item Hitung selang-$t$ $90$ persen untuk dua belas data waktu tunda pada bab ini,
        lalu bandingkan dengan selang $95$ persen. Berapa banyak data yang diperlukan
        agar lebar selang $95$ persen turun menjadi separuh, bila $s$ tidak berubah?
\end{enumerate}

Petunjuk untuk tiga butir tersulit. Butir 3: tulis
$\hat\theta_n-\theta=\operatorname{bias}_n+(\hat\theta_n-\mathbb{E}\hat\theta_n)$,
lalu pakai ketaksamaan segitiga dan Chebyshev pada suku kedua. Butir 6: ketaksamaan
pertama adalah $|\hat p-p_0|\le z\sqrt{\hat p(1-\hat p)/n}$, dan uji yang memakai
$\sqrt{p_0(1-p_0)/n}$ memakai penyebut berbeda, sehingga himpunan $p_0$ yang tidak
ditolak bukan lagi selang Wald melainkan himpunan yang lebih rumit. Butir 5: daerah
tolaknya adalah $|k-15|>1{,}96\sqrt{30\cdot0{,}25}$, jadi $k\le9$ atau $k\ge21$, dan
dayanya dihitung dengan menjumlahkan pmf binomial pada $p=0{,}7$ di kedua ekor itu.
\end{latihan}

\begin{untukpeneliti}
Yang benar-benar terbukti di bab ini: dekomposisi MSE; kekonsistenan dari MSE menuju
nol; cakupan tepat interval-$z$ pada model normal; bentuk tertutup interval Wilson dan
letaknya di $[0,1]$; aturan tiga sebagai batas eksak satu sisi beserta
ketaksamaan $p_{\mathrm{hi}}<\ln 20/n$; keseragaman nilai-p di bawah $H_0$ untuk
statistik kontinu; dan kesetaraan selang dengan uji pada pasangan yang disebutkan.
Yang hanya populer, bukan teorema: bahwa $95$ persen berarti peluang bersyarat pada
data yang sudah teramati; bahwa nilai-p besar mendukung $H_0$; dan bahwa lebar selang
yang lebih kecil selalu lebih baik, padahal perbandingan lebar tanpa cakupan tidak
berarti. Cakupan Wilson di bawah nominal pada beberapa $(n,p)$ juga fakta hitung, bukan
teorema.

Sumber primer: Wasserman, \emph{All of Statistics}, Bab 6 (Models, Statistical
Inference, and Learning) untuk bias, MSE, dan tafsiran selang; Bab 9 (Parametric
Inference) untuk interval dan metode plug-in; Bab 10 (Hypothesis Testing and p-values)
untuk nilai-p, galat tipe I dan II, daya, serta hubungan uji dengan selang. Interval
Wilson berasal dari E.~B. Wilson, ``Probable Inference, the Law of Succession, and
Statistical Inference'' (1927); aturan tiga dari B.~D. Jovanovic dan P.~S. Levy,
``A Look at the Rule of Three'' (1997). Keduanya tidak dibahas Wasserman, jadi
rujukannya disebut langsung, bukan lewat buku teks perantara.

Batas keberlakuan: seluruh bab bekerja dengan sampel bebas dan berdistribusi sama,
satu parameter, dan aproksimasi normal atau binomial yang diketahui. Interval-$z$
hanya tepat pada model normal dengan $\sigma$ diketahui; di luar itu ia pernyataan
asimtotik, dan pada $n$ kecil dengan distribusi berekor berat cakupannya dapat jauh di
bawah nominal. Interval Wilson adalah aproksimasi skor, bukan selang eksak; selang
eksak untuk proporsi dihitung dari sebaran binomial dan lebih lebar. Semua cakupan yang
dilaporkan bersyarat pada model binomial yang benar.

Pertanyaan terbuka yang layak dikerjakan pembaca: seberapa baik interval proporsi
berperilaku ketika $n$ bukan bilangan tetap melainkan berhenti secara adaptif, seperti
pada pengambilan sampel berurutan; dan bagaimana menyusun selang untuk metrik
klasifikasi yang dihitung dari data dependen, misalnya ketika satu contoh muncul di
beberapa lipatan validasi silang. Keduanya masih menjadi bahan penelitian.
\end{untukpeneliti}

\peta{Bab 20 pada Wasserman, \emph{All of Statistics}: Bab 6 (Models, Statistical
Inference, and Learning: penaksir, bias, MSE, galat baku, tafsiran selang), Bab 9
(Parametric Inference: interval kepercayaan dan metode plug-in), Bab 10 (Hypothesis
Testing and p-values: nilai-p, galat tipe I dan II, daya, kesetaraan uji dan selang).
Interval Wilson dan aturan tiga diambil dari makalah aslinya, Wilson (1927) dan
Jovanovic--Levy (1997), karena Wasserman tidak membahas keduanya. Hukum bilangan besar,
teorema limit pusat, dan ketaksamaan Chebyshev yang dipakai di sini dibuktikan di
Bab 18; peluang bersyarat dan kekeliruan pembalik syarat dibahas di Bab 16. Bab 25
melanjutkan bab ini untuk metrik klasifikasi, termasuk aturan tiga pada pelaporan
evaluasi model.}
